% =============================================================================
% APPENDIX VS: LIT-SMITH: Vernon Smith Research Integration
% =============================================================================
% Module: BCM2_LIT_SMITH (Literature Integration)
% Category: LIT
% Version: 1.0 (January 2026)
% Status: Final
% Language: English
%
% This appendix integrates Vernon Smith's foundational contributions to
% experimental economics, market design, and the two-rationalities framework
% into the Evidence-Based Framework (EBF).
%
% =============================================================================

% =============================================================================
% CROSS-REFERENCE STRUCTURE
% =============================================================================
%
% CHAPTER LINKAGE (Required):
% - Primary Chapter: Chapter 4 (Empirical Foundations)
% - Secondary Chapters: Chapter 5 (Complementarity), Chapter 9 (Context),
%                       Chapter 11 (Awareness), Chapter 14 (Behavioral Change)
%
% DEPENDENCIES (what this appendix REQUIRES):
% - Appendix K (LIT-FEHR): Social preferences experiments
% - Appendix U (LIT-KT): Decision-making foundations
% - Appendix UNMAPPED_RO (LIT-ROTH): Market design foundations
%
% DEPENDENTS (what REQUIRES this appendix):
% - Appendix UNMAPPED_LLM (METHOD-LLMMC): Experimental validation
% - Appendix E (METHOD-OPS): Laboratory methodology
% - Appendix Y (DOMAIN-CAPITAL): Asset market experiments
%
% =============================================================================

\begin{tcolorbox}[colback=purple!5!white, colframe=purple!75!black,
    title=Chapter Linkage]

\textbf{Primary Chapter:} Chapter 4: Empirical Foundations
\begin{itemize}[nosep]
\item Section 4.2: Experimental Evidence $\leftrightarrow$ This Appendix Section VS.2
\item Section 4.4: Laboratory Methods $\leftrightarrow$ This Appendix Section VS.3
\end{itemize}

\textbf{Secondary Chapters:}
\begin{itemize}[nosep]
\item Chapter 5: Complementarity --- market institution design
\item Chapter 9: Context --- ecological vs. constructivist rationality
\item Chapter 11: Awareness --- two-rationality framework
\item Chapter 14: Behavioral Change --- experimental methodology
\end{itemize}

\textbf{Reading Guide:}
\begin{itemize}[nosep]
\item For experimental methodology: Read this appendix first
\item For market design applications: See Appendix UNMAPPED_RO (LIT-ROTH)
\item For decision theory foundations: See Appendix U (LIT-KT)
\end{itemize}

\end{tcolorbox}

\chapter{LIT-SMITH: Vernon Smith Research Integration}
\label{app:lit-smith}


\begin{tcolorbox}[colback=blue!5!white, colframe=blue!75!black,
    title=Quick Reference: Appendix UNMAPPED_VS]

\textbf{Appendix:} VS (LIT)

\textbf{Core Question:} What are the key findings in this domain?

\textbf{Key Concepts:}
\begin{itemize}[nosep]
\item Framework integration
\item Parameter validation
\item Cross-references
\end{itemize}

\textbf{Cross-References:} See related appendices below
\end{tcolorbox}

\begin{abstract}
Vernon Smith (Nobel Prize 2002) pioneered experimental economics as a rigorous
methodology for testing economic theory. This appendix integrates his 60+ years
of research on market experiments, induced value theory, asset bubbles, and the
two-rationalities framework (constructivist vs. ecological) into EBF. Smith's
work provides the methodological foundation for laboratory validation of behavioral
predictions and demonstrates how market institutions generate emergent efficiency
even when individual rationality is limited.
\end{abstract}

\begin{tcolorbox}[colback=blue!5!white, colframe=blue!75!black,
    title=Quick Reference: LIT-SMITH]

\textbf{Core Question:} How can laboratory experiments test economic theory, and
what do market experiments reveal about the relationship between individual
cognition and institutional outcomes?

\textbf{Key Formula (Induced Value Theory):}
\begin{equation}
V_i(x) = R_i(x) - C_i \quad \text{where } R_i \text{ is experimenter-induced redemption value}
\label{eq:vs-induced-value}
\end{equation}

\textbf{Key Concepts:}
\begin{itemize}[nosep]
\item \textbf{Induced Value Theory}: Experimental control via monetary incentives
\item \textbf{Double Auction}: Efficient price discovery mechanism
\item \textbf{Two Rationalities}: Constructivist (deliberate) vs. Ecological (emergent)
\item \textbf{Market Efficiency}: Convergence despite bounded individual rationality
\end{itemize}

\textbf{Cross-References:} LIT-FEHR (K), LIT-KT (U), LIT-ROTH (RO), LIT-CAMERER (CA), METHOD-EXP
\end{tcolorbox}

% -----------------------------------------------------------------------------
% APPENDIX SCOPE BOX
% -----------------------------------------------------------------------------

\begin{tcolorbox}[
    colback=orange!5!white,
    colframe=orange!75!black,
    title={\textbf{Appendix Scope: Ziel / In-Scope / Out-of-Scope / Constraints}},
    fonttitle=\bfseries
]

\textbf{Ziel (Objective):}
\begin{itemize}[nosep]
    \item How do market institutions transform boundedly rational individual
behavior into efficient aggregate outcomes, and what are the conditions under
which this transformation fails?
\end{itemize}

\textbf{In-Scope:}
\begin{itemize}[nosep]
    \item The Fundamental Question
    \item Core Theory: The Smith Research Program
    \item Axioms: The Smith Framework
    \item Integration with EBF
    \item Worked Example
\end{itemize}

\textbf{Out-of-Scope (delegiert an andere Appendices):}
\begin{itemize}[nosep]
    \item LIT-FEHR $\rightarrow$ Appendix K
    \item LIT-KT $\rightarrow$ Appendix U
    \item LIT-ROTH $\rightarrow$ Appendix UNMAPPED_RO
    \item METHOD-LLMMC $\rightarrow$ Appendix UNMAPPED_LLM
    \item METHOD-OPS $\rightarrow$ Appendix E
\end{itemize}

\textbf{Constraints (Anwendungsgrenzen):}
\begin{itemize}[nosep]
    \item [TODO: Define when this appendix does NOT apply]
    \item [TODO: Define required prerequisites]
    \item [TODO: Define known limitations]
\end{itemize}

\textbf{Lieferobjekte (Deliverables):}
\begin{itemize}[nosep]
    \item 4 Table(s)
    \item 17 Equation(s)
    \item 6 Axiom(s)
    \item Worked Example(s)
\end{itemize}

\end{tcolorbox}




%%%%%%%%%%%%%%%%%%%%%%%%%%%%%%%%%%%%%%%%%%%%%%%%%%%%%%%%%%%%%%%%%%%%%%%%%%%%%%%
\section{The Fundamental Question}
\label{sec:vs-fundamental}
%%%%%%%%%%%%%%%%%%%%%%%%%%%%%%%%%%%%%%%%%%%%%%%%%%%%%%%%%%%%%%%%%%%%%%%%%%%%%%%

\subsection{Why This Matters}

Before Vernon Smith's pioneering work in the 1960s, economics was largely a
theoretical discipline with limited ability to test its predictions under
controlled conditions. Smith established that economic theories could be
rigorously tested in laboratory settings, creating a new empirical foundation
for the field.

His work matters for EBF because:
\begin{enumerate}
\item \textbf{Methodological Foundation}: Provides the experimental toolkit for
      validating behavioral predictions
\item \textbf{Institution-Behavior Interface}: Shows how market institutions
      shape individual behavior and aggregate outcomes
\item \textbf{Two-Rationalities Framework}: Reconciles behavioral anomalies
      with market efficiency through ecological rationality
\item \textbf{Bubble Dynamics}: Demonstrates systematic departures from
      fundamental value in asset markets
\end{enumerate}

\subsection{The Core Problem}

Classical economics assumed rational agents but had limited tools to test this
assumption. Behavioral economics documented systematic deviations from rationality.
Smith's contribution was showing that:

\begin{quote}
\textit{``Markets can be efficient even when individuals are not, because
institutions encode the experience and wisdom of past generations.''}
--- Smith (2003 Nobel Lecture)
\end{quote}

\begin{tcolorbox}[colback=red!5!white, colframe=red!75!black,
    title=The Fundamental Question]
\textbf{How do market institutions transform boundedly rational individual
behavior into efficient aggregate outcomes, and what are the conditions under
which this transformation fails?}
\end{tcolorbox}


%%%%%%%%%%%%%%%%%%%%%%%%%%%%%%%%%%%%%%%%%%%%%%%%%%%%%%%%%%%%%%%%%%%%%%%%%%%%%%%
\section{Core Theory: The Smith Research Program}
\label{sec:vs-theory}
%%%%%%%%%%%%%%%%%%%%%%%%%%%%%%%%%%%%%%%%%%%%%%%%%%%%%%%%%%%%%%%%%%%%%%%%%%%%%%%

\subsection{Induced Value Theory (1976)}

Smith's methodological breakthrough was \textit{induced value theory}---the
principle that experimenter-controlled rewards can create well-defined
preferences in laboratory subjects.

\textbf{The Three Precepts:}
\begin{enumerate}
\item \textbf{Monotonicity}: More reward medium (money) is preferred to less
\item \textbf{Salience}: Rewards depend on actions in the experiment
\item \textbf{Dominance}: Rewards dominate other motivations in the experiment
\end{enumerate}

\begin{equation}
U_i^{exp}(x, m) = U_i^{lab}(m | R_i(x)) \quad \text{where } R_i(x) \text{ is redemption schedule}
\label{eq:vs-utility}
\end{equation}

This methodology enables clean tests of market theory by controlling the
``demand and supply curves'' in the laboratory.

\subsection{Market Experiments: The Double Auction (1962)}

Smith's 1962 \textit{Journal of Political Economy} paper demonstrated that
double auction markets converge rapidly to competitive equilibrium, even with
small numbers of traders and incomplete information.

\textbf{Key Finding:} Markets with 4-6 traders on each side converge to within
5\% of competitive equilibrium within 3-4 trading periods.

\begin{equation}
|P_{observed} - P^*| < 0.05 \cdot P^* \quad \text{after } t \geq 3 \text{ periods}
\label{eq:vs-convergence}
\end{equation}

\textbf{Implication for EBF:} Market institutions can generate efficient outcomes
even when individual traders have limited information and bounded rationality.
The $\gamma$-complementarities in EBF operate through institutional channels.

\subsection{Asset Market Bubbles (1988)}

Smith, Suchanek, and Williams (1988) demonstrated that experimental asset
markets systematically generate price bubbles that deviate from fundamental
value:

\begin{equation}
P_t = P^*_t + B_t \quad \text{where } B_t = \text{bubble component}
\label{eq:vs-bubble}
\end{equation}

\textbf{Key Findings:}
\begin{itemize}
\item Bubbles persist even with experienced traders
\item Common knowledge of fundamentals does not eliminate bubbles
\item Bubbles are reduced but not eliminated by short-selling
\item Experience across multiple markets reduces bubble magnitude
\end{itemize}

\textbf{EBF Integration:} Bubbles represent a failure of $A(\cdot)$ (awareness)
at the collective level---traders are aware of fundamental value but do not
act on this knowledge due to speculative motives.

\subsection{Two Rationalities: Constructivist vs. Ecological (2003)}

Smith's Nobel lecture introduced the distinction between two forms of rationality:

\begin{table}[htbp]
\centering
\caption{Smith's Two Rationalities Framework}
\label{tab:vs-rationalities}
\renewcommand{\arraystretch}{1.4}
\begin{tabular}{@{}lll@{}}
\toprule
\textbf{Dimension} & \textbf{Constructivist} & \textbf{Ecological} \\
\midrule
Origin & Deliberate design & Emergent evolution \\
Knowledge & Explicit, articulable & Tacit, embedded in rules \\
Process & Top-down reasoning & Bottom-up adaptation \\
Examples & Planning, optimization & Markets, social norms \\
Failures & Complexity, hubris & Slow adaptation \\
EBF Connection & CORE-AWARE ($A$) & CORE-WHEN ($\Psi$) \\
\bottomrule
\end{tabular}
\end{table}

\begin{equation}
\text{Outcome} = f(\text{Constructivist}_i) \times g(\text{Ecological}_{institution})
\label{eq:vs-two-rationalities}
\end{equation}

\textbf{Key Insight:} Ecological rationality explains why markets work better
than individuals---institutions embody the accumulated wisdom of many
generations of trial-and-error learning.


%%%%%%%%%%%%%%%%%%%%%%%%%%%%%%%%%%%%%%%%%%%%%%%%%%%%%%%%%%%%%%%%%%%%%%%%%%%%%%%
\section{Axioms: The Smith Framework}
\label{sec:vs-axioms}
%%%%%%%%%%%%%%%%%%%%%%%%%%%%%%%%%%%%%%%%%%%%%%%%%%%%%%%%%%%%%%%%%%%%%%%%%%%%%%%

\begin{tcolorbox}[colback=gray!5!white, colframe=gray!75!black,
    title=Axiom UNMAPPED_VS-1: Induced Value Sufficiency]
\textbf{Statement:} Under monotonicity, salience, and dominance, experimenter-
induced values create well-defined preferences that enable clean tests of
economic theory.

\begin{equation}
\text{If } \frac{\partial U}{\partial m} > 0, \frac{\partial m}{\partial x} \neq 0, U^{lab} \gg U^{other} \Rightarrow V_i(x) = R_i(x) - C_i
\end{equation}

\textbf{Interpretation:} Laboratory experiments can achieve internal validity
comparable to natural sciences through proper incentive design.

\textbf{Implication:} EBF predictions can be tested in laboratory settings
using induced value methods.
\end{tcolorbox}

\begin{tcolorbox}[colback=gray!5!white, colframe=gray!75!black,
    title=Axiom UNMAPPED_VS-2: Institutional Efficiency]
\textbf{Statement:} Well-designed market institutions generate efficient
outcomes even when individual participants have bounded rationality.

\begin{equation}
\lim_{t \to \infty} |P_t - P^*| = 0 \quad \text{for double auction markets}
\end{equation}

\textbf{Interpretation:} Markets are ``smart'' even when traders are not---
efficiency is an emergent property of the institution, not a requirement of
individual rationality.

\textbf{Implication:} EBF's $\gamma$-complementarities operate through
institutional channels, not just individual cognition.
\end{tcolorbox}

\begin{tcolorbox}[colback=gray!5!white, colframe=gray!75!black,
    title=Axiom UNMAPPED_VS-3: Bubble Persistence]
\textbf{Statement:} Asset markets systematically generate price bubbles even
under conditions where fundamental value is common knowledge.

\begin{equation}
E[B_t | \text{common knowledge of } P^*] > 0
\end{equation}

\textbf{Interpretation:} Speculative dynamics create deviations from fundamental
value that are robust to information interventions.

\textbf{Implication:} CORE-AWARE ($A$) at collective level differs from
individual awareness---group dynamics amplify speculative motives.
\end{tcolorbox}

\begin{tcolorbox}[colback=gray!5!white, colframe=gray!75!black,
    title=Axiom UNMAPPED_VS-4: Ecological Rationality]
\textbf{Statement:} Social institutions encode ecological rationality that
compensates for individual cognitive limitations.

\begin{equation}
\text{Outcome}_{social} = f(\text{Individual}_{bounded}) \times g(\text{Institution}_{evolved})
\end{equation}

\textbf{Interpretation:} Evolution and learning at the institutional level
substitute for optimization at the individual level.

\textbf{Implication:} EBF's CORE-WHEN ($\Psi$) captures institutional context
that shapes behavior beyond individual cognition.
\end{tcolorbox}

\begin{tcolorbox}[colback=gray!5!white, colframe=gray!75!black,
    title=Axiom UNMAPPED_VS-5: Experience Learning]
\textbf{Statement:} Repeated market participation reduces behavioral anomalies
and improves market efficiency.

\begin{equation}
|B_t^{(n)}| < |B_t^{(n-1)}| \quad \text{where } n = \text{market experience}
\end{equation}

\textbf{Interpretation:} Traders learn from market feedback, and this learning
transfers across similar market environments.

\textbf{Implication:} CORE-STAGE (BCJ) dynamics apply to market behavior---
experience moves agents through stages of market sophistication.
\end{tcolorbox}

\begin{tcolorbox}[colback=gray!5!white, colframe=gray!75!black,
    title=Axiom UNMAPPED_VS-6: Humanomics]
\textbf{Statement:} Economic behavior is grounded in both self-interest (Wealth
of Nations) and social sentiments (Theory of Moral Sentiments), following Adam
Smith's dual framework.

\begin{equation}
U_i = U^{self}(x_i) + \alpha \cdot U^{social}(x_{-i}, \text{norms})
\end{equation}

\textbf{Interpretation:} Complete economic theory requires integrating both
self-interested and other-regarding motivations, as Adam Smith recognized.

\textbf{Implication:} EBF's INU/KNU/IDN decomposition formalizes this duality
in the welfare dimensions.
\end{tcolorbox}


%%%%%%%%%%%%%%%%%%%%%%%%%%%%%%%%%%%%%%%%%%%%%%%%%%%%%%%%%%%%%%%%%%%%%%%%%%%%%%%

%%%%%%%%%%%%%%%%%%%%%%%%%%%%%%%%%%%%%%%%%%%%%%%%%%%%%%%%%%%%%%%%%%%%%%%%%%%%%%%
\section{Key Results}
\label{sec:results}
%%%%%%%%%%%%%%%%%%%%%%%%%%%%%%%%%%%%%%%%%%%%%%%%%%%%%%%%%%%%%%%%%%%%%%%%%%%%%%%

The analysis yields the following key results:

\begin{enumerate}
    \item \textbf{Result 1:} [TODO: Describe main finding]
    \item \textbf{Result 2:} [TODO: Describe secondary finding]
    \item \textbf{Result 3:} [TODO: Describe implication]
\end{enumerate}

\section{Integration with EBF}
\label{sec:vs-integration}
%%%%%%%%%%%%%%%%%%%%%%%%%%%%%%%%%%%%%%%%%%%%%%%%%%%%%%%%%%%%%%%%%%%%%%%%%%%%%%%

\subsection{Connection to CORE-AWARE (AU)}

Smith's two-rationalities framework directly maps to EBF's awareness structure:

\begin{table}[htbp]
\centering
\caption{Smith Rationalities $\leftrightarrow$ EBF Awareness}
\label{tab:vs-aware}
\renewcommand{\arraystretch}{1.3}
\begin{tabular}{@{}lll@{}}
\toprule
\textbf{Smith Concept} & \textbf{EBF Component} & \textbf{Formalization} \\
\midrule
Constructivist Rationality & Explicit Awareness & $A^{explicit}(t^*)$ \\
Ecological Rationality & Implicit Awareness & $A^{implicit}(\Psi)$ \\
Market Intelligence & Collective Awareness & $A^{market} = \prod_i A_i$ \\
\bottomrule
\end{tabular}
\end{table}

\subsection{Connection to CORE-WHEN (V)}

Smith's institutional emphasis aligns with EBF's context framework:

\begin{equation}
\Psi_{institution} = (\text{rules}, \text{norms}, \text{information}, \text{feedback})
\label{eq:vs-psi}
\end{equation}

Market institutions shape $\Psi$ dimensions:
\begin{itemize}
\item \textbf{Rules}: Trading protocols (double auction, posted offer)
\item \textbf{Norms}: Expected behavior, fair pricing
\item \textbf{Information}: Price signals, order book transparency
\item \textbf{Feedback}: Profit/loss, market clearing
\end{itemize}

\subsection{Connection to CORE-HOW (B)}

Smith's market experiments reveal $\gamma$-complementarities:

\begin{equation}
\gamma_{buyer,seller} > 0 \quad \text{through price discovery mechanism}
\label{eq:vs-gamma}
\end{equation}

The double auction creates positive complementarity between buyer and seller
strategies---competitive bidding benefits both sides through efficient price
discovery.

\begin{tcolorbox}[colback=yellow!5!white, colframe=yellow!75!black,
    title=Integration: LIT-SMITH $\leftrightarrow$ EBF CORE]
\textbf{What LIT-SMITH provides:}
\begin{itemize}[nosep]
\item Experimental methodology for testing EBF predictions
\item Two-rationalities framework for understanding awareness
\item Institutional design principles for $\gamma$-optimization
\item Bubble dynamics showing collective awareness failures
\end{itemize}

\textbf{What it requires:}
\begin{itemize}[nosep]
\item CORE-AWARE for individual vs. collective awareness
\item CORE-WHEN for institutional context effects
\item CORE-HOW for complementarity through institutions
\item BBB for calibrated parameters from experiments
\end{itemize}
\end{tcolorbox}


%%%%%%%%%%%%%%%%%%%%%%%%%%%%%%%%%%%%%%%%%%%%%%%%%%%%%%%%%%%%%%%%%%%%%%%%%%%%%%%
\section{Worked Example}
\label{sec:vs-example}
%%%%%%%%%%%%%%%%%%%%%%%%%%%%%%%%%%%%%%%%%%%%%%%%%%%%%%%%%%%%%%%%%%%%%%%%%%%%%%%

\subsection{Asset Market Bubble and EBF Awareness}

\paragraph{Setup.} Consider an experimental asset market with 12 traders, where
the fundamental value of the asset declines linearly from \$3.60 to \$0.24 over
15 periods (dividend of \$0.24 per period).

\paragraph{Application.} Using Smith's bubble framework:

\textbf{Step 1: Individual Awareness}
\begin{equation}
A_i^{fundamental} = 0.8 \quad \text{(traders know dividend structure)}
\end{equation}

\textbf{Step 2: Collective Awareness Failure}
\begin{equation}
A^{collective} = 0.3 \quad \text{(speculative motives dominate)}
\end{equation}

\textbf{Step 3: Bubble Magnitude}
\begin{equation}
B_t = P_t - P^*_t = (1 - A^{collective}/A_i) \times P^*_t \approx 0.6 \times P^*_t
\end{equation}

\textbf{Step 4: Experience Effect}
After 2-3 market replications:
\begin{equation}
A^{collective}_{n=3} \approx 0.7 \Rightarrow B_t \approx 0.1 \times P^*_t
\end{equation}

\paragraph{Result.} The EBF awareness framework explains why bubbles persist
(collective awareness $<$ individual awareness) and why experience reduces them
(learning increases collective awareness).

\paragraph{Discussion.} Smith's experiments show that CORE-AWARE must distinguish
between individual and collective awareness levels. Market failures occur when
$A^{collective} \ll A^{individual}$, even when individuals understand fundamentals.


%%%%%%%%%%%%%%%%%%%%%%%%%%%%%%%%%%%%%%%%%%%%%%%%%%%%%%%%%%%%%%%%%%%%%%%%%%%%%%%
\section{Implications for Practice}
\label{sec:vs-practice}
%%%%%%%%%%%%%%%%%%%%%%%%%%%%%%%%%%%%%%%%%%%%%%%%%%%%%%%%%%%%%%%%%%%%%%%%%%%%%%%

\begin{tcolorbox}[colback=green!5!white, colframe=green!75!black,
    title=Practical Implications]
\begin{enumerate}
\item \textbf{Laboratory Testing}: EBF predictions should be validated through
      induced-value experiments before field implementation
\item \textbf{Institution Design}: Focus on institutional rules that generate
      ecological rationality, not just individual incentives
\item \textbf{Bubble Prevention}: Design interventions that increase collective
      awareness, not just individual information
\item \textbf{Experience Matters}: Allow for learning periods when introducing
      new market mechanisms or behavioral interventions
\item \textbf{Dual Motivation}: Recognize both self-interest and social
      sentiments in intervention design (humanomics)
\end{enumerate}
\end{tcolorbox}


%%%%%%%%%%%%%%%%%%%%%%%%%%%%%%%%%%%%%%%%%%%%%%%%%%%%%%%%%%%%%%%%%%%%%%%%%%%%%%%

% === AUTO-GENERATED ENTRIES (2026-02-22) ===

%%%%%%%%%%%%%%%%%%%%%%%%%%%%%%%%%%%%%%%%%%%%%%%%%%%%%%%%%%%%%%%%%%%%%%%%%%%%%%%
\subsubsection{1759: The Theory of Moral Sentiments}
%%%%%%%%%%%%%%%%%%%%%%%%%%%%%%%%%%%%%%%%%%%%%%%%%%%%%%%%%%%%%%%%%%%%%%%%%%%%%%%

\paragraph{Citation.}
Smith, Adam (1759). \textit{The Theory of Moral Sentiments}. A. Millar.

\paragraph{10C Integration.}
\begin{itemize}[nosep]
  \item \textbf{Domain}: [To be specified]
  \item \textbf{use\_for}: LIT-SMITH
\end{itemize}


%%%%%%%%%%%%%%%%%%%%%%%%%%%%%%%%%%%%%%%%%%%%%%%%%%%%%%%%%%%%%%%%%%%%%%%%%%%%%%%
\subsubsection{1776: An Inquiry into the Nature and Causes of the Wealth of Natio...}
%%%%%%%%%%%%%%%%%%%%%%%%%%%%%%%%%%%%%%%%%%%%%%%%%%%%%%%%%%%%%%%%%%%%%%%%%%%%%%%

\paragraph{Citation.}
Smith, Adam (1776). \textit{An Inquiry into the Nature and Causes of the Wealth of Nations}. W. Strahan and T. Cadell.

\paragraph{10C Integration.}
\begin{itemize}[nosep]
  \item \textbf{Domain}: [To be specified]
  \item \textbf{use\_for}: LIT-SMITH
\end{itemize}


%%%%%%%%%%%%%%%%%%%%%%%%%%%%%%%%%%%%%%%%%%%%%%%%%%%%%%%%%%%%%%%%%%%%%%%%%%%%%%%
\subsubsection{1988: Bubbles, Crashes, and Endogenous Expectations in Experimenta...}
%%%%%%%%%%%%%%%%%%%%%%%%%%%%%%%%%%%%%%%%%%%%%%%%%%%%%%%%%%%%%%%%%%%%%%%%%%%%%%%

\paragraph{Citation.}
Smith, Vernon L. and Suchanek, Gerry L. and Williams, Arlington W. (1988). ``Bubbles, Crashes, and Endogenous Expectations in Experimental Spot Asset Markets.'' \textit{Econometrica}, 56, 1119--1151.

\paragraph{10C Integration.}
\begin{itemize}[nosep]
  \item \textbf{Domain}: [To be specified]
  \item \textbf{use\_for}: LIT-SMITH
\end{itemize}


%%%%%%%%%%%%%%%%%%%%%%%%%%%%%%%%%%%%%%%%%%%%%%%%%%%%%%%%%%%%%%%%%%%%%%%%%%%%%%%
\subsubsection{1962: An Experimental Study of Competitive Market Behavior}
%%%%%%%%%%%%%%%%%%%%%%%%%%%%%%%%%%%%%%%%%%%%%%%%%%%%%%%%%%%%%%%%%%%%%%%%%%%%%%%

\paragraph{Citation.}
Smith, Vernon L. (1962). ``An Experimental Study of Competitive Market Behavior.'' \textit{Journal of Political Economy}, 70.

\paragraph{10C Integration.}
\begin{itemize}[nosep]
  \item \textbf{Domain}: [To be specified]
  \item \textbf{use\_for}: LIT-SMITH
\end{itemize}


%%%%%%%%%%%%%%%%%%%%%%%%%%%%%%%%%%%%%%%%%%%%%%%%%%%%%%%%%%%%%%%%%%%%%%%%%%%%%%%
\subsubsection{1965: Experimental Auction Markets and the Walrasian Hypothesis}
%%%%%%%%%%%%%%%%%%%%%%%%%%%%%%%%%%%%%%%%%%%%%%%%%%%%%%%%%%%%%%%%%%%%%%%%%%%%%%%

\paragraph{Citation.}
Smith, Vernon L. (1965). ``Experimental Auction Markets and the Walrasian Hypothesis.'' \textit{Journal of Political Economy}, 73.

\paragraph{10C Integration.}
\begin{itemize}[nosep]
  \item \textbf{Domain}: [To be specified]
  \item \textbf{use\_for}: LIT-SMITH
\end{itemize}


%%%%%%%%%%%%%%%%%%%%%%%%%%%%%%%%%%%%%%%%%%%%%%%%%%%%%%%%%%%%%%%%%%%%%%%%%%%%%%%
\subsubsection{1967: Experimental Studies of Discrimination versus Competition in...}
%%%%%%%%%%%%%%%%%%%%%%%%%%%%%%%%%%%%%%%%%%%%%%%%%%%%%%%%%%%%%%%%%%%%%%%%%%%%%%%

\paragraph{Citation.}
Smith, Vernon L. (1967). ``Experimental Studies of Discrimination versus Competition in Sealed-Bid Auction Markets.'' \textit{Journal of Business}, 40.

\paragraph{10C Integration.}
\begin{itemize}[nosep]
  \item \textbf{Domain}: [To be specified]
  \item \textbf{use\_for}: LIT-SMITH
\end{itemize}


%%%%%%%%%%%%%%%%%%%%%%%%%%%%%%%%%%%%%%%%%%%%%%%%%%%%%%%%%%%%%%%%%%%%%%%%%%%%%%%
\subsubsection{1976: Experimental Economics: Induced Value Theory}
%%%%%%%%%%%%%%%%%%%%%%%%%%%%%%%%%%%%%%%%%%%%%%%%%%%%%%%%%%%%%%%%%%%%%%%%%%%%%%%

\paragraph{Citation.}
Smith, Vernon L. (1976). ``Experimental Economics: Induced Value Theory.'' \textit{American Economic Review}, 66.

\paragraph{10C Integration.}
\begin{itemize}[nosep]
  \item \textbf{Domain}: [To be specified]
  \item \textbf{use\_for}: LIT-SMITH
\end{itemize}


%%%%%%%%%%%%%%%%%%%%%%%%%%%%%%%%%%%%%%%%%%%%%%%%%%%%%%%%%%%%%%%%%%%%%%%%%%%%%%%
\subsubsection{1980: Relevance of Laboratory Experiments to Testing Resource Allo...}
%%%%%%%%%%%%%%%%%%%%%%%%%%%%%%%%%%%%%%%%%%%%%%%%%%%%%%%%%%%%%%%%%%%%%%%%%%%%%%%

\paragraph{Citation.}
Smith, Vernon L. (1980). ``Relevance of Laboratory Experiments to Testing Resource Allocation Theory.'' \textit{Evaluation of Econometric Models}.

\paragraph{10C Integration.}
\begin{itemize}[nosep]
  \item \textbf{Domain}: [To be specified]
  \item \textbf{use\_for}: LIT-SMITH
\end{itemize}


%%%%%%%%%%%%%%%%%%%%%%%%%%%%%%%%%%%%%%%%%%%%%%%%%%%%%%%%%%%%%%%%%%%%%%%%%%%%%%%
\subsubsection{1982: Competitive Market Institutions: Double Auctions vs. Sealed ...}
%%%%%%%%%%%%%%%%%%%%%%%%%%%%%%%%%%%%%%%%%%%%%%%%%%%%%%%%%%%%%%%%%%%%%%%%%%%%%%%

\paragraph{Citation.}
Smith, Vernon L. and Williams, Arlington W. and Bratton, W. Kenneth and Vannoni, Michael G. (1982). ``Competitive Market Institutions: Double Auctions vs. Sealed Bid-Offer Auctions.'' \textit{American Economic Review}, 72.

\paragraph{10C Integration.}
\begin{itemize}[nosep]
  \item \textbf{Domain}: [To be specified]
  \item \textbf{use\_for}: LIT-SMITH
\end{itemize}


%%%%%%%%%%%%%%%%%%%%%%%%%%%%%%%%%%%%%%%%%%%%%%%%%%%%%%%%%%%%%%%%%%%%%%%%%%%%%%%
\subsubsection{1982: Microeconomic Systems as an Experimental Science}
%%%%%%%%%%%%%%%%%%%%%%%%%%%%%%%%%%%%%%%%%%%%%%%%%%%%%%%%%%%%%%%%%%%%%%%%%%%%%%%

\paragraph{Citation.}
Smith, Vernon L. (1982). ``Microeconomic Systems as an Experimental Science.'' \textit{American Economic Review}, 72.

\paragraph{10C Integration.}
\begin{itemize}[nosep]
  \item \textbf{Domain}: [To be specified]
  \item \textbf{use\_for}: LIT-SMITH
\end{itemize}


%%%%%%%%%%%%%%%%%%%%%%%%%%%%%%%%%%%%%%%%%%%%%%%%%%%%%%%%%%%%%%%%%%%%%%%%%%%%%%%
\subsubsection{1989: Theory, Experiment and Economics}
%%%%%%%%%%%%%%%%%%%%%%%%%%%%%%%%%%%%%%%%%%%%%%%%%%%%%%%%%%%%%%%%%%%%%%%%%%%%%%%

\paragraph{Citation.}
Smith, Vernon L. (1989). ``Theory, Experiment and Economics.'' \textit{Journal of Economic Perspectives}, 3.

\paragraph{10C Integration.}
\begin{itemize}[nosep]
  \item \textbf{Domain}: [To be specified]
  \item \textbf{use\_for}: LIT-SMITH
\end{itemize}


%%%%%%%%%%%%%%%%%%%%%%%%%%%%%%%%%%%%%%%%%%%%%%%%%%%%%%%%%%%%%%%%%%%%%%%%%%%%%%%
\subsubsection{1991: Rational Choice: The Contrast between Economics and Psycholo...}
%%%%%%%%%%%%%%%%%%%%%%%%%%%%%%%%%%%%%%%%%%%%%%%%%%%%%%%%%%%%%%%%%%%%%%%%%%%%%%%

\paragraph{Citation.}
Smith, Vernon L. (1991). ``Rational Choice: The Contrast between Economics and Psychology.'' \textit{Journal of Political Economy}, 99.

\paragraph{10C Integration.}
\begin{itemize}[nosep]
  \item \textbf{Domain}: [To be specified]
  \item \textbf{use\_for}: LIT-SMITH
\end{itemize}


%%%%%%%%%%%%%%%%%%%%%%%%%%%%%%%%%%%%%%%%%%%%%%%%%%%%%%%%%%%%%%%%%%%%%%%%%%%%%%%
\subsubsection{1994: Economics in the Laboratory}
%%%%%%%%%%%%%%%%%%%%%%%%%%%%%%%%%%%%%%%%%%%%%%%%%%%%%%%%%%%%%%%%%%%%%%%%%%%%%%%

\paragraph{Citation.}
Smith, Vernon L. (1994). ``Economics in the Laboratory.'' \textit{Journal of Economic Perspectives}, 8.

\paragraph{10C Integration.}
\begin{itemize}[nosep]
  \item \textbf{Domain}: [To be specified]
  \item \textbf{use\_for}: LIT-SMITH
\end{itemize}


%%%%%%%%%%%%%%%%%%%%%%%%%%%%%%%%%%%%%%%%%%%%%%%%%%%%%%%%%%%%%%%%%%%%%%%%%%%%%%%
\subsubsection{1994: Preferences, Property Rights, and Anonymity in Bargaining Ga...}
%%%%%%%%%%%%%%%%%%%%%%%%%%%%%%%%%%%%%%%%%%%%%%%%%%%%%%%%%%%%%%%%%%%%%%%%%%%%%%%

\paragraph{Citation.}
Smith, Vernon L. (1994). ``Preferences, Property Rights, and Anonymity in Bargaining Games.'' \textit{Games and Economic Behavior}, 7.

\paragraph{10C Integration.}
\begin{itemize}[nosep]
  \item \textbf{Domain}: [To be specified]
  \item \textbf{use\_for}: LIT-SMITH
\end{itemize}


%%%%%%%%%%%%%%%%%%%%%%%%%%%%%%%%%%%%%%%%%%%%%%%%%%%%%%%%%%%%%%%%%%%%%%%%%%%%%%%
\subsubsection{1998: The Two Faces of Adam Smith}
%%%%%%%%%%%%%%%%%%%%%%%%%%%%%%%%%%%%%%%%%%%%%%%%%%%%%%%%%%%%%%%%%%%%%%%%%%%%%%%

\paragraph{Citation.}
Smith, Vernon L. (1998). ``The Two Faces of Adam Smith.'' \textit{Southern Economic Journal}, 65.

\paragraph{10C Integration.}
\begin{itemize}[nosep]
  \item \textbf{Domain}: [To be specified]
  \item \textbf{use\_for}: LIT-SMITH
\end{itemize}


%%%%%%%%%%%%%%%%%%%%%%%%%%%%%%%%%%%%%%%%%%%%%%%%%%%%%%%%%%%%%%%%%%%%%%%%%%%%%%%
\subsubsection{2000: Bargaining and Market Behavior: Essays in Experimental Econo...}
%%%%%%%%%%%%%%%%%%%%%%%%%%%%%%%%%%%%%%%%%%%%%%%%%%%%%%%%%%%%%%%%%%%%%%%%%%%%%%%

\paragraph{Citation.}
Smith, Vernon L. and Wilson, Bart (2000). \textit{Bargaining and Market Behavior: Essays in Experimental Economics}. Publisher.

\paragraph{10C Integration.}
\begin{itemize}[nosep]
  \item \textbf{Domain}: [To be specified]
  \item \textbf{use\_for}: LIT-SMITH
\end{itemize}


%%%%%%%%%%%%%%%%%%%%%%%%%%%%%%%%%%%%%%%%%%%%%%%%%%%%%%%%%%%%%%%%%%%%%%%%%%%%%%%
\subsubsection{2002: Method in Experiment: Rhetoric and Reality}
%%%%%%%%%%%%%%%%%%%%%%%%%%%%%%%%%%%%%%%%%%%%%%%%%%%%%%%%%%%%%%%%%%%%%%%%%%%%%%%

\paragraph{Citation.}
Smith, Vernon L. (2002). ``Method in Experiment: Rhetoric and Reality.'' \textit{Experimental Economics}, 5.

\paragraph{10C Integration.}
\begin{itemize}[nosep]
  \item \textbf{Domain}: [To be specified]
  \item \textbf{use\_for}: LIT-SMITH
\end{itemize}


%%%%%%%%%%%%%%%%%%%%%%%%%%%%%%%%%%%%%%%%%%%%%%%%%%%%%%%%%%%%%%%%%%%%%%%%%%%%%%%
\subsubsection{2003: Constructivist and Ecological Rationality in Economics}
%%%%%%%%%%%%%%%%%%%%%%%%%%%%%%%%%%%%%%%%%%%%%%%%%%%%%%%%%%%%%%%%%%%%%%%%%%%%%%%

\paragraph{Citation.}
Smith, Vernon L. (2003). ``Constructivist and Ecological Rationality in Economics.'' \textit{American Economic Review}, 93.

\paragraph{10C Integration.}
\begin{itemize}[nosep]
  \item \textbf{Domain}: [To be specified]
  \item \textbf{use\_for}: LIT-SMITH
\end{itemize}


%%%%%%%%%%%%%%%%%%%%%%%%%%%%%%%%%%%%%%%%%%%%%%%%%%%%%%%%%%%%%%%%%%%%%%%%%%%%%%%
\subsubsection{2005: Hayek and Experimental Economics}
%%%%%%%%%%%%%%%%%%%%%%%%%%%%%%%%%%%%%%%%%%%%%%%%%%%%%%%%%%%%%%%%%%%%%%%%%%%%%%%

\paragraph{Citation.}
Smith, Vernon L. (2005). ``Hayek and Experimental Economics.'' \textit{Review of Austrian Economics}, 18.

\paragraph{10C Integration.}
\begin{itemize}[nosep]
  \item \textbf{Domain}: [To be specified]
  \item \textbf{use\_for}: LIT-SMITH
\end{itemize}


%%%%%%%%%%%%%%%%%%%%%%%%%%%%%%%%%%%%%%%%%%%%%%%%%%%%%%%%%%%%%%%%%%%%%%%%%%%%%%%
\subsubsection{2008: Rationality in Economics: Constructivist and Ecological Form...}
%%%%%%%%%%%%%%%%%%%%%%%%%%%%%%%%%%%%%%%%%%%%%%%%%%%%%%%%%%%%%%%%%%%%%%%%%%%%%%%

\paragraph{Citation.}
Smith, Vernon L. (2008). \textit{Rationality in Economics: Constructivist and Ecological Forms}. Publisher.

\paragraph{10C Integration.}
\begin{itemize}[nosep]
  \item \textbf{Domain}: [To be specified]
  \item \textbf{use\_for}: LIT-SMITH
\end{itemize}


%%%%%%%%%%%%%%%%%%%%%%%%%%%%%%%%%%%%%%%%%%%%%%%%%%%%%%%%%%%%%%%%%%%%%%%%%%%%%%%
\subsubsection{2010: Discovering Classical Economics}
%%%%%%%%%%%%%%%%%%%%%%%%%%%%%%%%%%%%%%%%%%%%%%%%%%%%%%%%%%%%%%%%%%%%%%%%%%%%%%%

\paragraph{Citation.}
Smith, Vernon L. (2010). ``Discovering Classical Economics.'' \textit{Journal of Economic Behavior and Organization}, 73.

\paragraph{10C Integration.}
\begin{itemize}[nosep]
  \item \textbf{Domain}: [To be specified]
  \item \textbf{use\_for}: LIT-SMITH
\end{itemize}


%%%%%%%%%%%%%%%%%%%%%%%%%%%%%%%%%%%%%%%%%%%%%%%%%%%%%%%%%%%%%%%%%%%%%%%%%%%%%%%
\subsubsection{2014: Adam Smith on Humanomic Behavior}
%%%%%%%%%%%%%%%%%%%%%%%%%%%%%%%%%%%%%%%%%%%%%%%%%%%%%%%%%%%%%%%%%%%%%%%%%%%%%%%

\paragraph{Citation.}
Smith, Vernon L. (2014). ``Adam Smith on Humanomic Behavior.'' \textit{Journal of Behavioral Finance}, 15.

\paragraph{10C Integration.}
\begin{itemize}[nosep]
  \item \textbf{Domain}: [To be specified]
  \item \textbf{use\_for}: LIT-SMITH
\end{itemize}


%%%%%%%%%%%%%%%%%%%%%%%%%%%%%%%%%%%%%%%%%%%%%%%%%%%%%%%%%%%%%%%%%%%%%%%%%%%%%%%
\subsubsection{2019: Humanomics: Moral Sentiments and the Wealth of Nations for t...}
%%%%%%%%%%%%%%%%%%%%%%%%%%%%%%%%%%%%%%%%%%%%%%%%%%%%%%%%%%%%%%%%%%%%%%%%%%%%%%%

\paragraph{Citation.}
Smith, Vernon L. and Wilson, Bart J. (2019). \textit{Humanomics: Moral Sentiments and the Wealth of Nations for the Twenty-First Century}. Publisher.

\paragraph{10C Integration.}
\begin{itemize}[nosep]
  \item \textbf{Domain}: [To be specified]
  \item \textbf{use\_for}: LIT-SMITH
\end{itemize}


%%%%%%%%%%%%%%%%%%%%%%%%%%%%%%%%%%%%%%%%%%%%%%%%%%%%%%%%%%%%%%%%%%%%%%%%%%%%%%%
\subsubsection{1999: Reconciling Psychology with Economics: Harnessing the Wind}
%%%%%%%%%%%%%%%%%%%%%%%%%%%%%%%%%%%%%%%%%%%%%%%%%%%%%%%%%%%%%%%%%%%%%%%%%%%%%%%

\paragraph{Citation.}
Smith, J.A. (1999). ``Reconciling Psychology with Economics: Harnessing the Wind.'' \textit{Journal}.

\paragraph{10C Integration.}
\begin{itemize}[nosep]
  \item \textbf{Domain}: [To be specified]
  \item \textbf{use\_for}: LIT-SMITH
\end{itemize}

% === END AUTO-GENERATED ENTRIES ===
\section{Summary}
\label{sec:vs-summary}
%%%%%%%%%%%%%%%%%%%%%%%%%%%%%%%%%%%%%%%%%%%%%%%%%%%%%%%%%%%%%%%%%%%%%%%%%%%%%%%

\begin{tcolorbox}[colback=green!10!white, colframe=green!75!black,
    title=\textbf{Appendix UNMAPPED_VS: Summary}]

\textbf{Core Question:} How do market institutions transform individual behavior
into aggregate outcomes?

\textbf{Main Contribution:}
\begin{itemize}[nosep]
    \item Established experimental economics as rigorous methodology
    \item Demonstrated market efficiency emerges from institutional design
    \item Identified bubble dynamics in asset markets
    \item Developed two-rationalities framework (constructivist vs. ecological)
    \item Integrated Adam Smith's dual view of human motivation (humanomics)
\end{itemize}

\textbf{Key Parameters:}
\begin{itemize}[nosep]
    \item Market convergence: within 5\% of equilibrium after 3-4 periods
    \item Bubble magnitude: 40-100\% above fundamentals in inexperienced markets
    \item Experience effect: 50-70\% bubble reduction after 2-3 replications
\end{itemize}

\textbf{Integration:} This appendix connects to CORE-AWARE (two rationalities),
CORE-WHEN (institutional context), CORE-HOW (market complementarities), and
BBB (experimental parameter calibration).

\end{tcolorbox}


%%%%%%%%%%%%%%%%%%%%%%%%%%%%%%%%%%%%%%%%%%%%%%%%%%%%%%%%%%%%%%%%%%%%%%%%%%%%%%%
\section{Glossary of Symbols}
\label{sec:vs-glossary}
%%%%%%%%%%%%%%%%%%%%%%%%%%%%%%%%%%%%%%%%%%%%%%%%%%%%%%%%%%%%%%%%%%%%%%%%%%%%%%%

\begin{tcolorbox}[colback=blue!5!white, colframe=blue!75!black]
\textbf{Central Reference:} For complete symbol definitions, see
\textbf{Appendix G} (Glossary, \S\ref{app:glossary}).

This section lists only symbols introduced or specialized in this appendix.
\end{tcolorbox}

\begin{table}[htbp]
\centering
\caption{Symbols Introduced in Appendix UNMAPPED_VS}
\label{tab:vs-symbols}
\renewcommand{\arraystretch}{1.3}
\begin{tabular}{@{}clp{6cm}c@{}}
\toprule
\textbf{Symbol} & \textbf{Name} & \textbf{Definition} & \textbf{Also in G?} \\
\midrule
$R_i(x)$ & Redemption Value & Experimenter-induced payoff function & NEW \\
$V_i(x)$ & Induced Value & Net value from experimental action & NEW \\
$B_t$ & Bubble Component & Deviation from fundamental value & NEW \\
$P^*$ & Fundamental Value & Intrinsic asset value & \checkmark \\
$A^{collective}$ & Collective Awareness & Market-level awareness & NEW \\
\bottomrule
\end{tabular}
\end{table}


%%%%%%%%%%%%%%%%%%%%%%%%%%%%%%%%%%%%%%%%%%%%%%%%%%%%%%%%%%%%%%%%%%%%%%%%%%%%%%%
\section{Critical Foundations}
\label{sec:vs-foundations}
%%%%%%%%%%%%%%%%%%%%%%%%%%%%%%%%%%%%%%%%%%%%%%%%%%%%%%%%%%%%%%%%%%%%%%%%%%%%%%%

\subsection{Objection 1: External Validity}

\textbf{Objection:} Laboratory experiments with student subjects and artificial
rewards cannot generalize to real-world markets with professional traders and
large stakes.

\textbf{Response:} Smith's parallelism principle addresses this: if the
theoretical mechanism is correctly specified, laboratory results should
generalize to any setting where that mechanism operates. Empirically, laboratory
findings (e.g., double auction efficiency, bubble patterns) have been replicated
in field settings and with professional traders.

\subsection{Objection 2: Institutional Specificity}

\textbf{Objection:} Smith's efficiency results depend on specific institutional
designs (double auction) that may not apply to other market structures.

\textbf{Response:} This is precisely Smith's point---institutions matter.
Different institutions produce different outcomes, and this institutional
sensitivity is a feature, not a bug. EBF incorporates this through $\Psi$-context.


%%%%%%%%%%%%%%%%%%%%%%%%%%%%%%%%%%%%%%%%%%%%%%%%%%%%%%%%%%%%%%%%%%%%%%%%%%%%%%%
\section{Open Issues}
\label{sec:vs-open}
%%%%%%%%%%%%%%%%%%%%%%%%%%%%%%%%%%%%%%%%%%%%%%%%%%%%%%%%%%%%%%%%%%%%%%%%%%%%%%%

\begin{table}[htbp]
\centering
\caption{Research Roadmap for LIT-SMITH Integration}
\label{tab:vs-roadmap}
\renewcommand{\arraystretch}{1.3}
\begin{tabular}{@{}clcl@{}}
\toprule
\textbf{\#} & \textbf{Issue} & \textbf{Priority} & \textbf{Status} \\
\midrule
1 & Collective vs. individual awareness calibration & HIGH & OPEN \\
2 & Institutional $\gamma$ measurement & HIGH & PARTIAL \\
3 & Experience learning curve parameters & MEDIUM & OPEN \\
4 & Cross-cultural replication & MEDIUM & PARTIAL \\
\bottomrule
\end{tabular}
\end{table}


%%%%%%%%%%%%%%%%%%%%%%%%%%%%%%%%%%%%%%%%%%%%%%%%%%%%%%%%%%%%%%%%%%%%%%%%%%%%%%%
\section{References}
\label{sec:vs-references}
%%%%%%%%%%%%%%%%%%%%%%%%%%%%%%%%%%%%%%%%%%%%%%%%%%%%%%%%%%%%%%%%%%%%%%%%%%%%%%%

\begin{tcolorbox}[colback=blue!5!white, colframe=blue!75!black]
\textbf{Master Bibliography:} For complete reference details, see
\texttt{bcm2\_master\_references.bib} in the repository.
\nocite{bcm_master}
\end{tcolorbox}

\subsection{Key References for This Appendix}

\begin{itemize}[nosep]
    \item \textbf{Experimental Methodology:} Smith (1962, 1976, 1982, 1994)
    \item \textbf{Asset Bubbles:} Smith, Suchanek \& Williams (1988), Lei et al. (2001)
    \item \textbf{Two Rationalities:} Smith (2003, 2008)
    \item \textbf{Humanomics:} Smith \& Wilson (2019), Smith (1998, 2014)
    \item \textbf{Market Institutions:} Plott \& Smith (1978), Smith et al. (1982)
\end{itemize}


%%%%%%%%%%%%%%%%%%%%%%%%%%%%%%%%%%%%%%%%%%%%%%%%%%%%%%%%%%%%%%%%%%%%%%%%%%%%%%%
% CROSS-REFERENCE FOOTER
%%%%%%%%%%%%%%%%%%%%%%%%%%%%%%%%%%%%%%%%%%%%%%%%%%%%%%%%%%%%%%%%%%%%%%%%%%%%%%%

\vspace{1em}
\hrule
\vspace{0.5em}

\noindent\textit{Cross-references:}
Appendix G (Glossary),
Appendix K (LIT-FEHR),
Appendix U (LIT-KT),
Appendix UNMAPPED_RO (LIT-ROTH),
Appendix UNMAPPED_CA (LIT-CAMERER),
Appendix Y (DOMAIN-CAPITAL).

\noindent\textit{Master Bibliography:} \texttt{bcm2\_master\_references.bib}

% =============================================================================
% END OF APPENDIX VS
% =============================================================================
